\documentstyle[%


righttag%


,11pt%


,amssymb%


,russian%               remove in Eng.


,izvru%                 Set izven% in Eng.


% ,izvru-ks             for brief communications


]{amsart}





%       Math definitions





\newcommand{\End}{\operatorname{End}}


\newcommand{\const}{\operatorname{const}}


\newcommand{\Ran}{\operatorname{Ran}}


\newcommand{\tts}[1]{}





\theoremstyle{plain}


\theorembodyfont{\it}


\newtheorem{thmnonum}{\hskip\parindent Теорема}


\renewcommand{\thethmnonum}{}





%\hoffset=-15mm                  For Eng.





\setcounter{page}{75}


\god{1997}


\nomer{10~(425)} %                Remove in Eng.


%\nomer{Vol.\,41, No.\,10}        Set in Eng.


% \str{75-81}                    Set in Eng.


\udk{517.983}


\author{\it Е.Л.\,Ульянова}


\title{О спектральных свойствах \\


относительно конечномерных возмущений\\


самосопряженных операторов}





\date{}


% \pagestyle{myheadings}         Set in Eng.





\pagestyle{plain} %              Remove in Eng.


\begin{document}





% \thispagestyle{plain}          Set in Eng.





\maketitle





% Body text








Пусть $H$ --- бесконечномерное комплексное гильбертово пространство и


$A:D(A)\subset H\to H$


--- самосопряженный полуограниченный дискретный оператор (т.\,е. оператор


$R(\lambda,A)=(A-\lambda I)^{-1}$, $\lambda\in\rho(A)$, является


компактным оператором), собственные


значения которого $0<\lambda_1<\lambda_2<\dotsb<\lambda_n<\dotsb$,


соответствующие собственным векторам $e_1,e_2,\dotsc,e_n,\dotsc$,


простые.





С использованием метода подобных операторов \cite{Bas1}--\cite{Dun}


исследуются спектральные свойства возмущенных операторов вида


\begin{equation}


A-B, \quad Bx=\sum^N_{m=1}(Ax,a_m)b_m,\quad x\in D(A),


\end{equation}


где $a_m$, $b_m$, $m=1,\dotsc,N<\infty$, --- наборы векторов из


гильбертова пространства $H$. Такое


возмущение $B$ называется относительно конечномерным. Получены оценки


собственных значений и собственных векторов операторов вида (1).





В основе метода подобных операторов лежит понятие допустимой тройки


$(\frak A,J,\Gamma)$


(см. \cite{Bas2}). Здесь $\frak A$ --- пространство допустимых


возмущений,


которому принадлежит рассматриваемый оператор $B$,


и $J:\frak A\to\frak A$, $\Gamma:\frak A\to\End H$ --- два


линейных трансформатора (т.\,е. линейные операторы, действующие в


пространстве линейных операторов).





Пусть $n$ --- некоторое натуральное число, $P_nx=(x,e_n)e_n$ --- проектор


Рисса, соответствующий


собственному значению $\lambda_n$. В качестве пространства возмущений


возьмем


линейное многообразие операторов из банахова пространства


$L_A(H)$ линейных


операторов, подчиненных оператору $A$ и представимых в виде


\begin{equation}


X=X_0A,\quad X\in\sigma_2(H),


\end{equation}


где $\sigma_2(H)$ --- идеал операторов Гильберта-Шмидта, и таких, что


\begin{equation}


\|P_iXP_j\|\le \const\alpha_i\beta_j \; \forall i,j,


\end{equation}


ненулевые последовательности $\{\alpha_i\}$, $\{\beta_j\}$


неотрицательных чисел принадлежат


пространству $l_2$ и выбраны так, чтобы оператор $B$


удовлетворял условию


(3). В частности, условие (3) выполнено, если


\begin{equation}


\alpha_i=\biggl(\,\sum^N_{m=1}|(e_i,a_m)|^2\biggr)^{1/2},\qquad


\beta_j=\biggl(\,\sum^N_{m=1}|(e_i,b_m)|^2\biggr)^{1/2}.


\end{equation}


Линейное пространство $\frak A$ нормируем, положив


\begin{equation}


\|X\|_*=\inf\{c>0:\|P_iX_0P_j\|\le c\alpha_i\beta_j \;


\forall i,j\}.


\end{equation}





Определим линейные трансформаторы $\Gamma_n:\frak A\to\End H$ и


$J_n:\frak A\to\frak A$


\begin{gather}


J_nX=P_nXP_n+P^nXP^n,\quad X\in\frak A,\quad


P^n=I-P_n,\quad J=J_n,\\


\Gamma X=P_nXS_n-S_nXP_n,\quad X\in\frak A,\quad \Gamma=\Gamma_n,


\end{gather}


где $S_n\in\End H$ определяется на векторах $e_k$, $k\ge 1$,


соотношениями


\begin{equation}


S_ne_n=0,\quad S_ne_k=


\frac{1}{\lambda_n-\lambda_k}e_k,\quad k\ne n.


\end{equation}





\begin{lem}


Тройка $(\frak A,J,\Gamma)$ является допустимой для оператора $A$, если


выполнены условия


\begin{align}


\gamma_1(n)&=\biggl(\,


\sum^\infty\begin{Sb}j=1\\ j\ne n\end{Sb}


\alpha^2_n\beta^2_j


\frac{|\lambda_j|^2}{|\lambda_n-\lambda_j|^2}+


\sum^\infty\begin{Sb}j=1\\ j\ne n\end{Sb}


\alpha^2_j\beta^2_n


\frac{|\lambda_j|^2}{|\lambda_n-\lambda_j|^2}\biggr)^{1/2}


<\infty,\\


%


\gamma_2(n)&=\max\biggl\{\alpha_n\beta_n\sup_{j\ne n}


\frac{|\lambda_j|}{|\lambda_n-\lambda_j|},\;


\sum^\infty\begin{Sb}j=1\\ j\ne n\end{Sb}


\alpha_j\beta_j


\frac{|\lambda_j|^2}{|\lambda_n-\lambda_j|^2}\biggr\}<\infty.


\end{align}


\end{lem}





{\bf Доказательство.} Проверим аксиомы допустимой тройки для


$(\frak A,J,\Gamma)$. Полнота


пространства $\frak A$ следует из полноты пространства


$\sigma_2(H)$, непрерывности


нормы и неравенства


\begin{equation}


\|X_0\|^2_2=\sum^\infty_{i,j=1}\|P_iX_0P_j\|^2\le


\|X\|^2_*\sum^\infty_{i,j=1}\alpha^2_i\beta^2_j.


\end{equation}


Отсюда, поскольку


$$


\|Xx\|^2=\|XAx\|^2 \le \|X_0\|^2_2\|Ax\| \le


\|X_0\|^2_2(\|Ax\|+\|x\|)\quad \forall x\in D(A),


$$


следует непрерывное вложение $\frak A$ в $L_A(H)$.





Покажем теперь, что $\Gamma X\in\End(H)$ и оценим норму трансформатора


$\Gamma$. Для этого


отметим, что имеет место равенство


$\|\Gamma Xx\|^2=\|P_nXS_nx\|^2+\|S_nXP_nx\|^2$. Оценим каждое слагаемое


в правой


части этого равенства. Из равенства Парсеваля, учитывая соотношение (9) и


свойства операторов из $\frak A$, имеем


$$


\|P_nX_0S_nx\|^2=\biggl(\,


\sum_{j\ge 1,\,j\ne n}|(x,e_j)|(|\lambda_j|/|\lambda_n-\lambda_j|)


\|P_nX_0P_j\|\biggr)^2.


$$


Отсюда в силу неравенства Коши и определения нормы в $\frak A$ получаем


$$


\|P_nX_0S_n\|^2\le


\sum_{j\ge 1,\,j\ne n}\alpha^2_n\beta^2_j


(|\lambda_j|/|\lambda_n-\lambda_j|)^2\|X\|^2_*.


$$


Для второго слагаемого, поскольку $P_nx=(x,e_n)e_n$, в силу соотношений


(2), (3), (8) и равенства $P^n=I-P_n$ получаем


$$


\|S_nXP_nx\|^2 \le \sum_{j\ge 1,\,j\ne n}\alpha^2_j\beta^2_n


(|\lambda_j|/|\lambda_n-\lambda_j|)^2\|X\|^2_*,


$$


следовательно, $\Gamma X$ --- ограниченный оператор и справедлива оценка


\begin{equation}


\|\Gamma X\|\le \Gamma_1(n)\|X\|_*.


\end{equation}


Докажем теперь, что операторы $X\Gamma Y$ и $\Gamma YX$ принадлежат


пространству $\frak A$ и справедливо неравенство


\begin{equation}


\max\{\|X\Gamma Y\|_*,\|\Gamma YX\|_*\}\le


\gamma_2(n)\|X\|_*\|Y\|_*.


\end{equation}


Для этого покажем, что для каждого из этих операторов имеет место


неравенство (3) с константой $\Gamma_2(n)$. Поскольку


$X\Gamma Y=X_0A\Gamma Y_0A$, то, используя


формулу для


$\Gamma$, получаем соотношения


$P_jX_0A\Gamma Y_0P_i=P_jX_0AP_nY_0S_nP_i$, $i\ne n$;


$P_jX_0A \Gamma Y_0P_i=P_jX_0AS_nY_0P_n$, $i=n$. Оценим норму


каждого оператора


в правой части последнего равенства. Из соотношений (3), (8) получаем


$$


\|P_jX_0AP_nY_0S_nP_i\|\le\sup_{i\ne n}


(|\lambda_n|/|\lambda_n-\lambda_i|)\beta_n\alpha_n


\|X\|_*\|Y\|_*\alpha_j\beta_i.


$$


Аналогично, используя равенство Парсеваля, соотношения (3), (8), имеем


$$


\|P_jX_0AS_nY_0P_n\|\le\sum_{k\ge 1,\,k\ne n}


\alpha_k\beta_k(|\lambda_k|/|\lambda_n-\lambda_k|)^2


\|X\|_*\|Y\|_*\alpha_j\beta_n.


$$


Таким образом, получена оценка для нормы оператора $X\Gamma Y$.


Аналогично


получается оценка для нормы оператора $\Gamma YX$.





Включение $\Ran(\Gamma X)\subset D(A)$, $X\in\frak A$, следует из вида


оператора


$\Gamma X$. В силу соотношения (8) непосредственно подстановкой


проверяется, что


$$


A\Gamma X-\Gamma XA=P_nXS_n(\lambda_nI-A)+


(\lambda_nI-A)S_nXP_n=P_nXP+PXP_n=X-JX.\quad\square


$$





Итак, $(\frak A,J,\Gamma)$ --- допустимая тройка. Рассмотрим возмущенный


оператор


$A-B$ и выпишем условия его подобия оператору вида $A-JX$, где


$X\in\frak A$.


Оператор


преобразования подобия будем искать в виде $U=I+\Gamma X$. Учитывая, что


$J$ является проектором, получаем уравнение для $X\in\frak A$


\begin{equation}


X=B\Gamma X-\Gamma XJB-\Gamma XJ(B\Gamma X),\quad X\in\frak A.


\end{equation}


Пространство возмущений является прямой суммой подпространств вида


$\frak A_{ij}=\{Q_iXQ_j:X\in\frak A\}$,


$i,j=1,2$, $Q_1=P_n$, $Q_2=P^n$. Уравнение (14) перепишем для операторных


блоков $X_{ij}=Q_iXQ_j$, $i,j=1,2$, оператора $X\in\frak A$


\begin{align}


X_{11}&=B_{12}\Gamma X_{21}+B_{11},\\


X_{21}&=B_{22}\Gamma X_{21}-(\Gamma X_{21})B_{11}-


(\Gamma X_{21})B_{12}(\Gamma X_{21})+B_{21},\\


X_{22}&=B_{21}\Gamma X_{12}+B_{22},\\


X_{12}&=B_{11}\Gamma X_{12}-(\Gamma X_{12})B_{12}-


(\Gamma X_{12})B_{21}(\Gamma X_{12})+B_{21}.


\end{align}





Рассмотрим нелинейные операторы


$F_{ij}:\frak A_{ij}\to\frak A_{ij}$ ($i=2$, $j=1$ или


$i=1$, $j=2$), задаваемые


правыми частями уравнений (16) и (18) соответственно.





\begin{lem}


Пусть $n\ge 1$ --- некоторое натуральное число такое, что выполнено


условие


\begin{equation}


\varepsilon=\gamma(n)\|B\|_*<1/4,\quad


\gamma(n)=\max\{\gamma_1(n),\gamma_2(n)\}.


\end{equation}


Тогда оператор $A-B$, $B\in\frak A$, подобен оператору $A-JX^*$, оператор


преобразования подобия имеет вид


\begin{equation}


U=I+\Gamma X^*,\quad X^*=X^*_{11}+X^*_{12}+X^*_{21}+X^*_{22},


\end{equation}


$X^*$ --- решение уравнения {\em (14)}, а $X^*_{ij}$, $i,j=1,2$, ---


решение систем


{\em (16)--(18)}.


\end{lem}





\begin{pf}


Из проведенных рассуждений следует, что отображение


$F_{ij}:\frak A_{ij}\to\frak A_{ij}$ ($i=2$, $j=1$ или $i=1$,


$j=2$) является сжимающим в шаре


\begin{gather*}


S_{ij}=\{Y\in\frak A_{ij}:\|Y-B_{ij}\|_* \le r\|B_{ij}\|_*\}, \\


r=\varepsilon(4+2\varepsilon)/(1-2\varepsilon-2\varepsilon^2+


(1-4\varepsilon)^{1/2}),


\end{gather*}


с константой сжатия $q=1-(1-4\varepsilon)^{1/2}$ и переводит его в себя.


Поэтому по теореме Банаха


(\cite{Kant}, c.\,605) оно имеет в этом шаре единственную неподвижную


точку $X^*_{ij}$ (т.\,е. уравнения (17) и (19) имеют единственное


решение) и справедлива оценка


\begin{equation}


\|X^*_{ij}-F^k_{ij}(0)\|_* \le


(q^k/(1-q))\|B_{ij}\|_*,\quad k\ge 1,\quad


F_{ij}(0)=B_{ij}.


\end{equation}





Непрерывная обратимость оператора $U$ следует из неравенства (12) и


свойства решения уравнений (16) и (18).


\end{pf}





Приведем теперь оценки собственных значений и собственных векторов


оператора $A-B$.





\begin{thmnonum}


Пусть для некоторого натурального $n\ge 1$ выполнено условие {\em (19)}.


Тогда имеют место оценки


\begin{equation}


|\Bar\lambda_n-\lambda_n+(Be_n,e_n)|\le |\lambda_n|


\beta_n\alpha_n\gamma_2(n)\|B_{21}\|_*\|B_{12}\|_*/


(1-4\gamma_2(n)\|B\|_*)^{1/2},


\end{equation}


\begin{multline}


\biggl\| \Bar e_n-e_n+\sum_{j\ge 1,\,j\ne n}(Be_n,e_j)Se_j\biggr\| \le\\


%


\le 4|\lambda_n|\beta_n\gamma_2(n)\|B_{21}\|_*\|B\|_*/


(1-4\gamma_2(n)\|B\|_*)^{1/2}\biggl(\,


\sum_{j\ge 1,\,j\ne n}(|\lambda_j|/|\lambda_n-\lambda_j|)^2


\biggr)^{1/2}.


\end{multline}


\end{thmnonum}





\begin{pf}


Из подобия оператора $A-B$ оператору $A-JX^*$ следует


$\Bar e_n=(I-\Gamma X^*)e_n$ и


$\Bar\lambda_n=((A-X^*_{11})e_n,e_n)$. Отсюда,


учитывая вид решения $X^*$ и трансформатора $\Gamma$, получаем


\begin{gather}


|\Bar\lambda_n-\lambda_n+(Be_n,e_n)|\le


\|B_{12}\Gamma B_{21}e_n\|+\|B_{12}\Gamma(X^*_{21}-B_{21})e_n\|,\\


\|\Bar e_n-e_n+S_nB_{21}e_n\|=\|S_n(X^*_{21}-B_{21})e_n\|.


\end{gather}


Из неравенств (24), (13), (21) и вида константы сжатия $q$


следует оценка (22).





Используя равенство Парсеваля и соотношения (3), (8), (21), получаем


$$


\|S_n(X^*_{21}-B_{21})e_n\|\le


(\beta_n\lambda_nq/(1-q))\biggl(\,


\sum_{j\ge 1,\,j\ne n}(|\lambda_j|/|\lambda_n-\lambda_j|)^2


\biggr)^{1/2}\|B_{21}\|_*.


$$


Окончательная оценка следует из соотношения (25) и вида константы сжатия


$q$.


\end{pf}





\begin{thebibliography}{9}





\bibitem{Bas1}


Баскаков А.Г. {\em Гармонический анализ линейных операторов}. Учеб.


пособие. -- Воронеж: Изд-во Воронежск. ун-та, 1987. -- 164 с.


\bibitem{Bas2}


Баскаков А.Г. {\em Спектральный анализ возмущенных неквазианалитических и


спектральных операторов} // Изв. РАН. Сер. матем. -- 1994. -- Т.\,58. --


\No\,4. -- С.\,3--32.


\bibitem{Dun}


Данфорд Н., Шварц Дж.Т. {\em Линейные операторы}. Т.\,3. -- М.: Мир,


1974. -- 661~с.


\bibitem{Kant}


Канторович Л.В., Акилов Г.П. {\em Функциональный анализ}. -- М.: Наука,


1984. -- 752~с.





\end{thebibliography}





\vskip 2cm





\post{\it Воронежский государственный университет}{\it Поступила}


\post{}{02.02.1995}





\end{document}


